%&pdftex
% JUAN LUIS VARONA
% ENVIADO A CONSTRUCTIVE APPROXIMATION
% PRIMERA VEZ: MAYO DE 1991.
% VERSION CORREGIDA
% REENVIADO EN DICIEMBRE DE 1991.
% ACEPTADO EN MAYO DE 1992.


% ------------------------------------------------


\magnification=\magstep1
\overfullrule=0pt 

\font\msym=msbm10
\def\R{\hbox{\msym R}} % real numbers
\def\N{\hbox{\msym N}} % integer numbers
\def\supp{\mathop{\rm supp}} % support
\def\span{\mathop{\rm span}} % span
\def\Id{\mathop{\rm Id}} % identity map
\def\llave#1{\left\{{#1}\right\}} % to put between braces
\def\normlp#1#2#3{\Vert{#1}\Vert_{L^{#2}({#3})}} % norm in L^p of a function
\def\th{{\cal H}_\alpha} % Hankel transform of order \alpha
\def\section#1{\goodbreak{\bf#1}\medskip}


% ------------------------------------------------


\centerline{\bf FOURIER SERIES OF FUNCTIONS}
\centerline{\bf WHOSE HANKEL TRANSFORM IS SUPPORTED ON $[0,1]$}

\bigskip

\centerline{JUAN L. VARONA\footnote{*}{The author has been supported
in part by DGICYT Grant PB89-0181-C02-02.}}
\footnote{}{PAPER PUBLISHED IN: {\it Constr.\ Approx.}\ {\bf 10} (1994), 65--75.}

\bigskip


{\narrower\noindent{\bf Abstract.} Let $J_\mu$ denote the Bessel function of
order $\mu$. For $\alpha > -1$, the system
$x^{-\alpha/2-1/2}J_{\alpha+2n+1}(x^{1/2})$, $n = 0, 1, 2, \dots$ is orthogonal
on $L^2((0,\infty),x^{\alpha}\,dx)$. In this paper we study the mean
convergence of Fourier series with respect to this system for functions whose
Hankel transform is supported on $[0,1]$. \smallskip}


\bigskip

\noindent{\sl 1991 Mathematics Subject Classification}: Primary 42C10; Secondary
44A05. \medskip

\noindent{\sl Key words and phrases}: Bessel functions, Fourier
series, Hankel transform, $A_p$-theory.

\bigskip


% ------------------------------------------------


\section {\S 1. Introduction.}


Let $J_{\mu}(x)$ stand for the Bessel function of order $\mu$ and
$\{P^{(\alpha,\beta)}_n(x)\}_{n=0}^\infty$ for the Jacobi polynomials (see [15]
and Ch.~VII and X in [5]). It is well known that the Jacobi polynomials are
orthogonal on (-1,1) with respect to the weight $(1-x)^\alpha(1+x)^\beta$,
$\alpha,\beta>-1$, and the Bessel functions satisfy the orthogonality relation
$$
   \int_0^\infty  J_{\alpha +2n+1}(x) J_{\alpha +2m+1}(x)\,{dx \over x}
   = {{\delta _{nm}} \over {2(2n+\alpha +1)}},
   \quad n,m=0,1,2,\ldots \qquad (\alpha>-1).
$$
If we denote
$$
   j_n^\alpha(x)=\sqrt{\alpha+2n+1}J_{\alpha+2n+1}(\sqrt{x})x^{-\alpha/2-1/2},
   \quad n=0,1,2,\ldots \leqno(1)
$$
then the system $\{j_n^\alpha\}_{n=0}^\infty$ is orthonormal on $L^2 (x^\alpha)
= L^2((0,\infty),x^{\alpha}\,dx)$.

The Bessel functions and the Jacobi polynomials are related by the formula (see
[6])
$$
   \int_0^\infty J_{\alpha +2n+1}(t) J_\alpha (xt)\,dt =
   x^\alpha P_n^{(\alpha,0)}(1-2x^2) \chi_{[0,1]}(x).
$$

Following [4] we define the Hankel transform $\th$ of order $\alpha>-1$ to be
the integral operator
$$
   \th(f,x) = {{x^{-\alpha /2}} \over 2}
   \int_0^\infty {f(t)J_\alpha (\sqrt {xt})t^{\alpha /2}\,dt},
   \quad x > 0, \leqno(2)
$$
for suitable functions $f$.

This means that
$$
   \th(j_n^\alpha,x) = \sqrt{\alpha+2n+1} P_n^{(\alpha,0)}(1-2x)
   \chi_{[0,1]}(x), \leqno(3)
$$
and therefore  $\supp (\th(j_n^\alpha)) \subseteq [0,1]$.

We consider the partial sums of the Fourier series with respect to the system
$\{j_n^\alpha\}_{n=0}^\infty$:
$$
   S_n(f,x) = \sum_{k=0}^n {c_k(f) j_k^\alpha (x)}, \qquad
   c_k(f) = \int_0^\infty {f(t) j_k^\alpha(t) t^\alpha\,dt}.
$$
They can be written as
$$
   S_n(f,x) = \int_0^\infty {f(t) K_n(x,t) t^\alpha\,dt}, \quad
   {\rm where \quad}
   K_n(x,t) = \sum_{k=0}^n j_k^\alpha(x)j_k^\alpha(t).\leqno(4)
$$
Series of this kind are a particular case of series $\sum_{n\ge0} a_n
J_{\alpha+n}$, which are usually called Neumann series. A study of their
pointwise convergence can be found in [16] and [17].

The main aim in this paper is to study the convergence of $S_nf$ in the
$L^p(x^\alpha)$-norm. This involves two problems:

a) To obtain uniform boundedness of the operators $S_nf$ in $L^p(x^\alpha)$.

b) To find the subspace of $L^p(x^\alpha)$ consisting of the functions
$f$ which can be approximated in the $L^p(x^\alpha)$-norm by its Fourier
series, that is, to describe the space
$$
   B_{p,\alpha} = \overline{\span} \{j_n^\alpha(x)\}_{n=0}^\infty
   \quad ({\rm closure\ in\ } L^p(x^\alpha)).
$$

In order to solve a) the kernel $K_n$ is decomposed in a suitable way which
reduces the problem to show the boundedness of the Hilbert transform with
weights, and hence some estimates for the Bessel functions and some results on
$A_p$ theory are needed. Some of these ideas have been used in the literature
(see [1], [8], [9], [10], [12], [13]).

Regarding to b), looking at (3) we only need to deal with
functions with Hankel transform supported on $[0,1]$. This leads us to
consider, in a natural way, the analogous of the disc multiplier for the 
Hankel transform, i.~e., the operator $M_{\alpha}$ defined by
$$
   \th(M_\alpha f,x) = \th(f,x) \chi_{[0,1]}(x).
$$

Our problem of expanding a function whose Hankel transform is supported on
$[0,1]$ with respect to an orthogonal system is in some sense similar to
expanding a function whose Fourier transform is supported on $[-1,1]$, which has
been treated in [1] using as an orthogonal system the spherical Bessel functions
$\sqrt{\pi\over 2x}J_{n+1/2}(x)$. The method we use to establish our results is
easier than the one in [1]. Our development in \S 2 can be adapted to
simplify some of the proofs in [1] and [2].

The paper is organized as follows: In section \S 2, we solve problem a). In
section \S 3 we prove that the operator $M_\alpha$, defined for suitable
functions $f$, can be extended from $L^p(x^\alpha)$ into itself and it turns 
out to be the projection operator. This allows us to study problem b) obtaining,
in section \S 4, a characterization of $B_{p,\alpha}$ in terms of $M_\alpha$ and
solving the convergence of $S_nf$ to $f$ in the $L^p(x^\alpha)$-norm.


% ------------------------------------------------

\bigskip


\section {\S 2. Uniform boundedness of the partial sums.}

In what follows, $\alpha\geq -{1 \over 2}$ and $p_0 = { {4(\alpha+1)} \over
{2\alpha+3}}$, $p_1 = { {4(\alpha+1)} \over {2\alpha+1}}$. From the well
known estimates (see [5] or [15])
$$
   J_\mu(x) = {x^\mu \over 2^\mu \Gamma(\mu+1)} + O(x^{\mu+2}),
   \quad x\to 0+
$$
and
$$
   J_\mu(x) = \sqrt{2\over \pi x}
   \left[\cos\left(x-{\mu\pi\over 2}-{\pi\over 4}\right) + O(x^{-1})\right],
   \quad x\to\infty \leqno(5)
$$
it follows
$$
   |J_\alpha(x)|\le C_\alpha x^\alpha, \quad x\in (0,\infty) \leqno(6)
$$
and
$$
   |J_\alpha(x)|\le C_\alpha x^{-1/2}, \quad x\in (0,\infty). \leqno(7)
$$

Given $p \in (1,\infty)$ and a fixed interval $(a,b)$, a weight $w$ is said to
belong to the $A_p(a,b)$ class if
$$
   \left(\int_I w(x)\,dx \right) \left(\int_I w(x)^{-1/(p-1)}\,dx \right)^{p-1}
   \le C|I|^p
$$
for every interval $I \subseteq (a,b)$, with $C$ independent of $I$. An
important application of $A_p$ theory lies in its relation with the boundedness
of the Hilbert transform
$$
   H(f,x) = \int_a^b { f(t) \over x-t }\,dt.
$$
Indeed, in [11] (see also [7] for further information) it is proved that
$$
   H : L^p((a,b),w) \longrightarrow L^p((a,b),w) {\rm\ bounded} \quad
   \Longleftrightarrow \quad w \in A_p(a,b).
$$
Besides, the norm of the Hilbert transform operator and the constant in the
$A_p$ definition depend only one on each other. This allows us to use the uniform
$A_p$ theory in a similar way to [10]. Let us suppose that a family of weights
$\{w_n\}_{n=0}^\infty$ defined in the same interval $(a,b)$ satisfies the $A_p$
condition with the same constant $C$ (in this case we will say that $w_n \in
A_p(a,b)$ uniformly). Then, the Hilbert transform $H$ is uniformly bounded from
$L^p((a,b),w_n)$ into itself, that is, with constant independent of $n$.

It is well known
$$
   x^\beta \in A_p(0,1) \Longleftrightarrow x^\beta \in A_p(1,\infty)
   \Longleftrightarrow x^\beta \in A_p(0,\infty) \Longleftrightarrow
   -1< \beta < p-1. \leqno(8)
$$
Moreover, by making the change of variable $x=r_nz$ we easily obtain that, if
$r_n\searrow 0$, then
$$
   \left( {|x|+r_n} \right)^\beta \in A_p(-1,1) {\rm\ uniformly} \
   \Longleftrightarrow \ -1<\beta <p-1. \leqno(9)
$$
By using (8), it is easy to prove that, if $\alpha\ge-1/2$ and $p_0 < p <p_1$,
then
$$
   x^{\alpha-\alpha p/2+p/4} \in A_p(0,\infty) {\rm \quad and \quad}
   x^{\alpha -\alpha p/2-p/4}\in A_p(0,\infty).
$$

The following result will be used later.

\bigskip

\proclaim Lemma 1. Let $\alpha \ge -1/2$, $\max \{{4\over 3},p_0 \} < p <
\min\{ 4,p_1 \}$ and $0<s_n \nearrow \infty$. Then
$$
   x^{\alpha-\alpha p/2+p/8}
   \left( |x^{1/2}-s_n|+s_n^{1/3} \right)^{p/4} \in A_p(0,\infty)
   {\sl\ uniformly} \leqno(10)
$$
and
$$
   x^{\alpha-\alpha p/2-p/8}
   \left( |x^{1/2}-s_n|+s_n^{1/3} \right)^{-p/4} \in A_p(0,\infty)
   {\sl\ uniformly}. \leqno(11)
$$

{\it Proof}.

Let us prove (11); the proof of (10) is similar. Making the change
of variable $x=s_n^2z$ in the $A_p$ definition, it is easy to show that (11) is
equivalent to proving
$$
   z^{\alpha-\alpha p/2-p/8}\left( {|z^{1/2}-1|+s_n^{-2/3}} \right)^{-p/4}
   \in A_p(0,\infty ) {\rm\ uniformly}.
$$
Taking into account the behaviour of this expression on the intervals
$(0,1/2)$, $(1/2,3/2)$ and $(3/2,\infty)$, it is not difficult to check that we
only need to see
$$
   z^{\alpha-\alpha p/2-p/8}\in A_p(0,{\textstyle{1 \over 2}}),\ \ 
   \left({|z^{1/2}-1|+s_n^{-2/3}} \right)^{-p/4}
   \in A_p({\textstyle{1 \over 2}},{\textstyle{3 \over 2}}), \ \ 
   z^{\alpha-\alpha p/2-p/4}\in A_p({\textstyle{3 \over 2}},\infty).
$$
The first and the third ones are true by (8). Finally, the second one follows
from (9) by using $|z^{1/2}-1|\sim|z-1|$ and making a change of variable.

\bigskip

\proclaim Theorem 1. Let $1 < p < \infty$ and $\alpha \geq -1/2$. Then, there
exists a constant $C$ independent of $n$ and $f$, such that
$$
   \Vert S_nf \Vert_{L^p(x^\alpha)} \leq C \Vert f \Vert_{L^p(x^\alpha)}
   \qquad \forall f \in L^p(x^\alpha) \leqno(12)
$$
if and only if $\max\{{4 \over 3}, p_0\} < p < \min\{4,p_1\}$.

{\it Proof.}

To find necessary conditions for (12) we apply the standard
argument used for the first time in [12]. The uniform boundedness implies that
of the operator $T_n=S_n-S_{n-1}$,
$$
   T_n(f,x) = c_n(f)j_n^\alpha(x) =
   j_n^\alpha(x) \int_0^\infty f(t) j_n^\alpha(t) t^\alpha\,dt.
$$
Thus, by using duality we can easily obtain that (12) implies
$$
   \normlp{j_n^\alpha}{p}{x^\alpha} \normlp{j_n^\alpha}{q}{x^\alpha}
   \le C.\leqno(13)
$$
Taking $n=0$ and applying (5) and (1), we find that $p_0 < p < p_1$.
Now, provided that $p_0 < p < p_1$, asymptotic estimates for $J_{\alpha+2n+1}(x)
$
with $n$ and $x$ large enough allow us to show that
$$
   \normlp{j_n^\alpha}{p}{x^\alpha} \sim \left\{\matrix{
   n^{-1-\alpha +2\alpha /p+2/p}      \hfill &{\hbox{if}\ p<4}\hfill \cr
   n^{-1/2-\alpha /2}(\log n)^{1/4}   \hfill &{\hbox{if}\ p=4}\hfill \cr
   n^{-5/6-\alpha +2\alpha /p+4/(3p)} \hfill &{\hbox{if}\ p>4}\hfill \cr}
   \right..
$$
This, together with (13), implies $4/3<p<4$.

On the other hand, let us suppose that $\max\{{4 \over 3}, p_0\} < p <
\min\{4,p_1\}$ and prove the uniform boundedness of $S_nf$.

We need a suitable decomposition of the kernel $K_n(x,t)$. In [16]
it is proved that
$$
   \phi_\mu(x,t) = \sum_{k=0}^\infty 2(\mu+2k+1) J_{\mu+2k+1}(x)
   J_{\mu+2k+1}(t)
$$
satisfies
$$
   \phi_\mu(x,t) = { x t \over x^2-t^2 }
   \llave{ {x J_{\mu+1}(x) J_\mu(t) - t J_\mu(x) J_{\mu+1}(t)} }.
$$
Consequently, by (4) and (1) it is clear that
$$
   K_n(x,t) = {1\over2} x^{-\alpha/2-1/2} t^{-\alpha/2-1/2}
   \llave{ {\phi _\alpha (x^{1/2},t^{1/2})
   - \phi _{\alpha + 2n+2}(x^{1/2},t^{1/2})}}.
$$
Now, by using $\ z J_{\alpha+2n+3}(z) =
(\alpha+2n+2)J_{\alpha+2n+2}(z)-zJ'_{\alpha+2n+2}(z)$ we obtain
$$
   \phi_{\alpha+2n+2} (x,t) = { x t \over x^2-t^2}
   \llave{ {J_{\alpha+2n+2}(x) t J'_{\alpha+2n+2}(t)
   - J_{\alpha+2n+2}(t) x J'_{\alpha+2n+2}(x)}}.
$$
Hence it follows
$$
   K_n(x,t) = { x^{-\alpha/2} t^{-\alpha/2} \over 2(x-t) }
   \llave{ x^{1/2} J_{\alpha +1}(x^{1/2}) J_\alpha (t^{1/2})
   - t^{1/2} J_\alpha(x^{1/2}) J_{\alpha +1}(t^{1/2})}
$$
$$
   + { x^{-\alpha/2} t^{-\alpha/2} \over 2(x-t)}
   \llave{ x^{1/2} J'_{\alpha +2n+2}(x^{1/2}) J_{\alpha+2n+2}(t^{1/2})
   - t^{1/2} J_{\alpha+2n+2}(x^{1/2}) J'_{\alpha +2n+2}(t^{1/2})}.
$$
This enables us to write
$$
   S_n(f,x) = W_1(f,x) - W_2(f,x) + W_{3,n}(f,x) - W_{4,n}(f,x),
$$
where
$$
   W_1(f,x) = {1\over2} \int_0^\infty x^{-\alpha/2+1/2} t^{-\alpha/2}
   { J_{\alpha +1}(x^{1/2}) J_\alpha (t^{1/2}) \over x-t } f(t) t^\alpha\,dt,
$$
$$
   W_2(f,x) = {1\over2} \int_0^\infty x^{-\alpha/2} t^{-\alpha/2+1/2}
   { J_\alpha (x^{1/2}) J_{\alpha +1}(t^{1/2}) \over x-t } f(t) t^\alpha\,dt,
$$
$$
   W_{3,n}(f,x) = {1\over2} \int_0^\infty x^{-\alpha/2+1/2} t^{-\alpha/2}
   { J'_\nu (x^{1/2}) J_\nu (t^{1/2}) \over x-t } f(t) t^\alpha\,dt,
$$
$$
   W_{4,n}(f,x) = {1\over2} \int_0^\infty x^{-\alpha/2} t^{-\alpha/2+1/2}
   { J_\nu (x^{1/2}) J'_\nu (t^{1/2}) \over x-t } f(t) t^\alpha\,dt
$$
and $\nu=\alpha+2n+2$.

Therefore, to prove (12) it is sufficient to see
$$
   \normlp{W_if}{p}{x^\alpha} \le C \normlp{f}{p}{x^\alpha}, \qquad i = 1,2
$$
and
$$
   \normlp{W_{i,n}f}{p}{x^\alpha} \le C \normlp{f}{p}{x^\alpha},\qquad i = 3,4.
$$

In order to apply $A_p$ theory we will use the following estimates for
Bessel functions and their derivatives:
$$
   |J_\nu (x)| \le Cx^{-1/4} \left( {|x-\nu |+\nu ^{1/3}} \right)^{-1/4},
   \leqno(14)
$$
$$
   |J'_\nu (x)| \le Cx^{-3/4} \left( {|x-\nu |+\nu ^{1/3}} \right)^{1/4},
   \leqno(15)
$$
where, again, $\nu=\alpha+2n+2$ and the constant $C$ depends only on $\alpha$.
These inequalities can be easily deduced from those used in [1] (see also [15]).

Now, the boundedness of the operators $W_1$ and $W_2$ is equivalent to
$$
   \normlp{x^{1/2-\alpha /2} J_{\alpha +1}(x^{1/2})
   H(t^{\alpha /2}J_\alpha (t^{1/2})f(t),x)}
   {p}{x^\alpha}
   \le C \normlp{f(x)}{p}{x^\alpha}
$$
and
$$
   \normlp{x^{-\alpha /2} J_\alpha (x^{1/2})
   H(t^{1/2+\alpha /2} J_{\alpha +1}(t^{1/2})f(t),x)}
   {p}{x^\alpha}
   \le C \normlp{f(x)}{p}{x^\alpha}.
$$
By using (7) and $x^{\alpha-\alpha p/2-p/4}\in A_p(0,\infty)$ it follows
$$
   \normlp{x^{-\alpha /2}J_\alpha (x^{1/2})
   H(t^{1/2+\alpha /2}J_{\alpha +1}(t^{1/2})f(t),x)}{p}{x^\alpha}^p
$$
$$
   \le C_1 \int_0^\infty \left| x^{-\alpha/2-1/4}
   H(t^{1/2+\alpha /2} J_{\alpha+1} (t^{1/2}) f(t),x) \right|^p x^\alpha\,dx
$$
$$
   \le C_2 \int_0^\infty \left| x^{1/2+\alpha /2}
   J_{\alpha+1} (x^{1/2}) f(x) \right|^p x^{\alpha-\alpha p/2-p/4}\,dx
$$
$$
   \le C_3\int_0^\infty \left| x^{1/2+\alpha /2-1/4} f(x) \right|^p
   x^{\alpha-\alpha p/2-p/4}\,dx
   = C_3 \normlp{f(x)}{p}{x^\alpha}^p.
$$
The verification of the other inequality is similar by using (7) and
$x^{\alpha-\alpha p/2+p/4}\in A_p(0,\infty)$.

To finish the proof we will prove the uniform boundedness of $W_{4,n}$; that of
$W_{3,n}$ can be obtained in a similar way. Taking
$g(t)=J'_\nu(t^{1/2})t^{\alpha/2+1/2}f(t)$ and using consecutively (14), (11)
and (15) we have
$$
   \normlp{W_{4,n}f}{p}{x^\alpha}^p = 2^{-p} \int_0^\infty 
   \left| \int_0^\infty  x^{-\alpha/2} t^{\alpha/2+1/2}
   { J_\nu(x^{1/2}) J'_\nu(t^{1/2}) \over x-t}
   f(t)\,dt \right|^p x^\alpha\,dx
$$
$$
   = 2^{-p} \int_0^\infty
   |H(g,x)|^p \left| J_\nu(x^{1/2}) \right|^p x^{-\alpha p/2+\alpha}\,dx
$$
$$
   \le C_1 \int_0^\infty 
   |H(g,x)|^p x^{\alpha-\alpha p/2-p/8}
   \left( |x^{1/2}-\nu|+\nu^{1/3} \right)^{-p/4}\,dx
$$
$$
   \le C_2 \int_0^\infty |g(x)|^p x^{\alpha-\alpha p/2-p/8}
   \left( |x^{1/2}-\nu|+\nu^{1/3} \right)^{-p/4}\,dx
   \le C_3 \normlp{f}{p}{x^\alpha}^p
$$
and the result follows.

\bigskip


% ------------------------------------------------


\section {\S 3. The projection operator $M_\alpha$.}

By using (6), the integral operator (2) exists for every $f\in
L^1(x^\alpha)$ and $\th f\in L^\infty(x^\alpha)$. Actually, we have
$$
    \normlp{\th f}{\infty}{x^\alpha}\le \textstyle{1\over 2}C_\alpha
    \normlp{f}{1}{x^\alpha}.\leqno(16)
$$
However (2) does not exist for $f\in L^p(x^\alpha)$ in general and so we will
introduce the operator $M_\alpha$ in a similar way to that of the disc
multiplier.

We consider the space
$$
   S^+ = \llave
   {f\in C^{\infty}(0,\infty) : \forall k,n\ge0,\; |t^k f^{(n)}(t)|<C_{k,n}}
$$
with the topology generated by the seminorms $\|\cdot\|_{k,n},\ k, n \in \N$,
defined by $\|f\|_{k,n}=\sup_{t\in(0,\infty)}t^k|f^{(n)}(t)|$ (see [3] and [4]).
It is easy to identify $S^+$ with the functions $f$ such that $f(t)=\phi(t)$,
$t\ge0$, for some $\phi(t)$ in the Schwartz class $S$.

With this notation $\th$ is an isomorphism of $S^+$ onto itself and $\th^2$ is
the identity map. Moreover, Fubini's Theorem implies the multiplication formula
for the Hankel transform:
$$
   \int_0^\infty \th(f,x) g(x) x^\alpha\,dx =
   \int_0^\infty \th(g,x) f(x) x^\alpha\,dx, \quad f,g \in S^+.\leqno(17)
$$
Now, taking $g(x)=\th(f,x)$ and using that $\th^2 = \Id$, we have Parseval's
formula $\normlp{\th f}{2}{x^\alpha}=\normlp{f}{2}{x^\alpha}$.

Since $S^+$ is dense in $L^2(x^\alpha)$, the operator $\th:S^+ \rightarrow S^+$
can be extended to $\th : L^2(x^\alpha) \rightarrow L^2(x^\alpha)$ satisfying
$$
   \normlp{\th f}{2}{x^\alpha} = \normlp{f}{2}{x^\alpha}
   \quad {\rm and} \quad \th^2 = \Id.\leqno(18)
$$

Let $M_\alpha$ be the operator defined by
$$
   M_\alpha(f,x) = \th(\chi_{[0,1]}\th f,x), \quad f \in S^+.
$$
Note that if $f\in S^+$ then $\th f \in S^+$; thus $\chi_{[0,1]}\th f\in
L^2(x^\alpha)$ and $M_\alpha f$ is well defined $\forall f\in S^+$.

\bigskip

\proclaim Theorem 2. Let $\alpha \geq -1/2$ and $p_0 < p <p_1$. Then there
exists a constant $C_{p,\alpha}$ such that
$$
   \Vert M_\alpha f \Vert_{L^p(x^\alpha)} \leq
   C_{p,\alpha} \Vert f \Vert_{L^p(x^\alpha)}, \quad \forall f \in S^+.
$$
Therefore, $M_\alpha$ can be extended to an operator (also denoted $M_\alpha$)
bounded on $L^p(x^\alpha)$ such that
\medskip
\item{\rm i)} $\th(M_\alpha f)=\th(f)\chi_{[0,1]}$ for all $f\in L^2(x^\alpha)
\cap L^p(x^\alpha)$.
\medskip
\item{\rm ii)} $M_\alpha^2 f = M_\alpha f$ for all $f \in L^p(x^\alpha)$.
\medskip
\item{\rm iii)} Moreover, for $f\in L^p(x^\alpha)$ and $g\in L^q(x^\alpha)$,
$1/p+1/q=1$, we have
$$
    \int_0^\infty{f(x)M_\alpha (g,x)x^\alpha\,dx} =
    \int_0^\infty{g(x)M_\alpha (f,x)x^\alpha\,dx}.
$$

{\it Proof.}

By using (2) and Fubini's Theorem we obtain
$$
   M_\alpha(f,x) = {1 \over 4} \int_0^\infty
   \left( \int_0^1 J_\alpha (\sqrt{yt}) J_\alpha (\sqrt {yx})\,dy \right)
   t^{\alpha /2} x^{-\alpha /2} f(t)\,dt.
$$
Moreover, a change of variable in Lommel's formula
$$
   \eqalign{
   \int_0^1 J_\alpha(yt) J_\alpha(yx) y\,dy
   &= {1 \over t^2-x^2}
   \left( J_\alpha(t) x J'_\alpha(x) - J_\alpha (x) t J'_\alpha (t) \right) \cr
   &= {1 \over t^2-x^2}
   \left( t J_{\alpha +1}(t) J_\alpha (x) - x J_\alpha(t) J_{\alpha +1}(x)
   \right)
   \cr}
$$
(for the last equality, use $zJ'_\alpha(z)=\alpha
J_\alpha(z)-zJ_{\alpha+1}(z)$) leads us to
$$
   \eqalign{
   M_\alpha(f,x) &= {1 \over 2} \int_0^\infty {
   t^{1/2} J_{\alpha+1}(t^{1/2}) J_\alpha(x^{1/2})
   - x^{1/2} J_\alpha(t^{1/2}) J_{\alpha +1}(x^{1/2})
   \over t-x}
   t^{\alpha /2} x^{-\alpha /2} f(t)\,dt \cr
   &= W_1(f,x) - W_2(f,x),
   \cr}
$$
where $W_1$ and $W_2$ are bounded operators (see the proof of Theorem 1).

Now, taking into account that $\th^2 = \Id$ and using standard density
arguments, the statements i) and ii) follow easily. Let us now prove iii):

From (17), it easily follows that, for $f,g\in L^2(x^\alpha)$, we have
$$
    \int_0^\infty{\th(f,x)\th(g,x)x^\alpha\,dx} =
    \int_0^\infty{f(x)g(x)x^\alpha\,dx}.\leqno(19)
$$

Now, let $U_1$ and $U_2$ be the bilinear functionals on
$L^p(x^\alpha)\times L^q(x^\alpha)$ defined by
$$
    U_1(f,g)=\int_0^\infty{f(x)M_\alpha(g,x)x^\alpha\,dx}
$$
and
$$
    U_2(f,g)=\int_0^\infty{g(x)M_\alpha(f,x)x^\alpha\,dx}.
$$
It suffices to show that $U_1$ and $U_2$ are bounded and
coincide on the subset $\big(L^2(x^\alpha)\times L^2(x^\alpha)\big)\cap
\big(L^p(x^\alpha)\times L^q(x^\alpha)\big)$, which is dense in
$L^p(x^\alpha)\times L^q(x^\alpha)$.
The boundedness of $U_1$ and $U_2$ is clear by H\"older's inequality
and $p_0<p,q<p_1$. Furthermore, by using (19), i) twice, and (19) once again, we
have
$$
    \int_0^\infty{f(x)M_\alpha(g,x)x^\alpha\,dx} =
    \int_0^\infty{\th(f,x)\th(M_\alpha g,x)x^\alpha\,dx} =
    \int_0^1{\th(f,x)\th(g,x)x^\alpha\,dx}
$$
$$
    =\int_0^\infty{\th(M_\alpha f,x)\th(g,x)x^\alpha\,dx} =
    \int_0^\infty{M_\alpha(f,x)g(x)x^\alpha\,dx}
$$
and so the proof is complete.

\bigskip


% ------------------------------------------------


\section {\S 4. Main consequences.}

\proclaim Definition. Let $\alpha\ge -1/2$ and $p_0 < p < p_1$. We define
$$
E_{p,\alpha} = \llave{ f \in L^p(x^\alpha) : M_\alpha f = f}
$$
endowed with the topology induced by $L^p(x^\alpha)$.

\bigskip

\proclaim Proposition 1. Let $\alpha\ge -1/2$, $p_0 < s < r < p_1$. Then
$E_{s,\alpha}\subset E_{r,\alpha}$ and the inclusion is continuous and dense.

{\it Proof.}

From (16) and (18), by using interpolation (see [14]), it follows
$$
    \normlp{\th f}{q}{x^\alpha}\le C \normlp{f}{p}{x^\alpha},
    \quad 1\le p \le 2\le q \le \infty,
    \quad \textstyle{1\over p}+\textstyle{1\over q}=1.
$$
On the other hand, if $p\le2$ and $p_0<p<p_1$, for $f\in S^+$ we have
$$
    \normlp{M_\alpha f}{\infty}{x^\alpha} =
    \normlp{\th(\chi_{[0,1]}\th f)}{\infty}{x^\alpha}
$$
$$
    \le C_1 \normlp{\chi_{[0,1]}\th f}{1}{x^\alpha} \le C_2
    \normlp{\th f}{q}{x^\alpha} \le C_3 \normlp{f}{p}{x^\alpha}.
$$
By density,
$$
    \normlp{M_\alpha f}{\infty}{x^\alpha}
    \le C \normlp{f}{p}{x^\alpha}.
$$
Since $M_\alpha$ is bounded from $L^p(x^\alpha)$ into itself, $p_0<p<p_1$, by
interpolation we obtain
$\normlp{M_\alpha f}{r}{x^\alpha}\le C \normlp{f}{p}{x^\alpha}$, $r\ge p$.
Interpolating one more time, we find
$$
    \normlp{M_\alpha f}{r}{x^\alpha}\le C\normlp{f}{s}{x^\alpha},
    \quad  p_0 < s < r < p_1.\leqno(20)
$$
If $f\in E_{s,\alpha}$ then $M_\alpha f=f$ and so (20) leads to
$\normlp{f}{r}{x^\alpha}\le C \normlp{f}{s}{x^\alpha}$, which implies
$E_{s,\alpha}\subset E_{r,\alpha}$.

To prove that the inclusion is dense, let $f\in E_{r,\alpha}$. Since
$L^r(x^\alpha) \cap L^s(x^\alpha)$ is dense in $L^r(x^\alpha)$, for each
$\varepsilon>0$ there exists a function $g\in L^r(x^\alpha) \cap L^s(x^\alpha)$
such that $\normlp{f-g}{r}{x^\alpha} < \varepsilon$. Taking $h=M_\alpha g$ it
follows that $h\in E_{p,\alpha}$. Then
$$
    \normlp{f-h}{r}{x^\alpha} = \normlp{M_\alpha f-M_\alpha g}{r}{x^\alpha}
    \le \|M_\alpha\|\;\normlp{f-g}{r}{x^\alpha} < C\varepsilon
$$
and the proof is complete.

\bigskip

If $1/p+1/q=1$, the following result says that the dual space $(E_{p,\alpha})'$
is isomorphic to $E_{q,\alpha}$ in the standard sense

\bigskip

\proclaim Proposition 2. Let $\alpha\ge-1/2$ and $T$ a bounded linear operator
on $E_{p,\alpha}$, $p_0<p<p_1$. There exists a unique function $g\in
E_{q,\alpha}$, $1/p+1/q=1$, such that
$$
    T(f)=\int_0^\infty{f(x)g(x)x^\alpha\,dx}
    \quad\forall f\in E_{p,\alpha}.\leqno(21)
$$
Furthermore, $C\normlp{g}{q}{x^\alpha}\le \|T\| \le \normlp{g}{q}{x^\alpha}$ for
a constant $C>0$ depending only on the norm of $M_\alpha$ in $L^q(x^\alpha)$.

{\it Proof.}

Let $T\in(E_{p,\alpha})'$. By Hahn-Banach's Theorem, $T$ can be extended to
$T\in(L^p(x^\alpha))'$ preserving its norm $\|T\|$. By duality between
$L^p(x^\alpha)$ and $L^q(x^\alpha)$, there exists $h\in L^q(x^\alpha)$ such that
$$
    T(f)=\int_0^\infty{f(x)h(x)x^\alpha\,dx}\quad\forall f\in L^p(x^\alpha)
$$
and $\normlp{h}{q}{x^\alpha}=\|T\|$.

If $g=M_\alpha h$ from ii) in Theorem 2 it follows $g\in E_{q,\alpha}$.
We are going to check that this function satisfies our purposes. If $f\in
E_{p,\alpha}$, from iii) in Theorem 2 we have
$$
    T(f) = \int_0^\infty{f(x)h(x)x^\alpha\,dx} =
    \int_0^\infty{M_\alpha(f,x)h(x)x^\alpha\,dx}
$$
$$
    = \int_0^\infty{M_\alpha(h,x)f(x)x^\alpha\,dx} =
    \int_0^\infty{g(x)f(x)x^\alpha\,dx}.
$$
By using $p_0<q<p_1$ and H\"older's inequality we have
$\|T\|\le \normlp {M_\alpha h}{q}{x^\alpha}$.
Thus, the equivalence of norms $\normlp{g}{q}{x^\alpha} \sim \|T\|$ follows
immediately.

To prove the uniqueness, suppose $p\ge2$ and there exist $g$ and $g'$ in
$E_{q,\alpha}$ satisfying (21). Then
$$
    \int_0^\infty{f(x)(g(x)-g'(x))x^\alpha\,dx} = 0
    \quad \forall f\in E_{p,\alpha}.
$$
Taking $f=g-g'\in E_{q,\alpha}\subset E_{p,\alpha}$ it follows $g-g'=0$ a.~e.
Now, the case $p<2$ is a simple consequence of this considering that
$L^q(x^\alpha)$ is reflexive and, since $E_{q,\alpha}$ is closed, then
$E_{q,\alpha}$ is also reflexive.

\bigskip

As it was pointed out in Introduction, let $B_{p,\alpha}$ stands for the
closure in $L^p(x^\alpha)$ of the space $\span\llave{j_n^\alpha}_{n=0}^\infty$.
The following results show the main consequences in this section.

\bigskip

\proclaim Corollary. Let $\alpha\geq -1/2$. Then
$\lim_{n\to\infty}\normlp{S_nf-f}{p}{x^\alpha}=0$ for every $f\in B_{p,\alpha}$
if and only if $\max\llave{{4\over 3},p_0} < p < \min\llave{4,p_1}$.
Moreover, if it is the case, then $B_{p,\alpha} = E_{p,\alpha}$.

{\it Proof.}

i) By using (3), (18) and the completeness of the Jacobi system it follows that
$\llave{j_n^\alpha}_{n=0}^\infty$ is complete in $E_{2,\alpha}$ and so
$B_{2,\alpha}=E_{2,\alpha}$. Hence the convergence for $p=2$ is clear.
When $p_0<p<p_1$ then $j_n^\alpha\in E_{p,\alpha}$ and therefore
$B_{p,\alpha}\subset E_{p,\alpha}$.

If $2<p<p_1$ and $f\in E_{p,\alpha}$ it follows that for each $\varepsilon>0$
there exists a function $g\in L^2(x^\alpha)\cap L^p(x^\alpha)$ such that
$\normlp{f-g}{p}{x^\alpha}<\varepsilon$. Taking $h=M_\alpha g$ we have $M_\alpha
h = h$ and so $h\in E_{2,\alpha}\cap E_{p,\alpha} = B_{2,\alpha}\cap
E_{p,\alpha}$. Since $M_\alpha$ is continuous then $\normlp{f-h}{p}{x^\alpha} =
\normlp{M_\alpha f-M_\alpha g}{p}{x^\alpha} < C\varepsilon$.
As $h\in B_{2,\alpha}$ it follows that there exists
$h'\in \span\llave{j_n^\alpha}_{n=0}^\infty$ such that
$\normlp{h-h'}{2}{x^\alpha}<\varepsilon$. By applying triangle inequality and
Proposition 1 we obtain $\normlp{f-h'}{p}{x^\alpha}<C_1\varepsilon$ and
therefore $E_{p,\alpha}\subset B_{p,\alpha}$. Then the mean convergence for $2
< p < \min\llave{4,p_1}$ is an immediate consequence of Theorem 1 and
Banach-Steinhaus Theorem.

If $\max\llave{{4\over 3},p_0}<p<2$, the result is easily obtained by duality.


% ------------------------------------------------

\bigskip

\goodbreak

\centerline {\bf REFERENCES}

\frenchspacing

\bigskip

\item {{\bf [1]}} J. A. Barcel\'o and A. C\'ordoba, {\it Band-limited
functions: $L^p$-convergence}, Trans. Amer. Math. Soc. {\bf 313} (1989),
655--669.

\item {{\bf [2]}} A. C\'ordoba, {\it The disc multiplier}, Duke Math. J. {\bf
58} (1989), 21--27.

\item {{\bf [3]}} A. J. Dur\'an, {\it Laguerre expansions of tempered
distributions and generalized functions}, J. Math. Anal. Appl. (to
appear).

\item {{\bf [4]}} A. J. Dur\'an, {\it On Hankel transform}, Proc. Amer. Math.
Soc. {\bf 110} (1990), 417--424.

\item {{\bf [5]}} A. Erd\'elyi et al., {\it ``Higher Transcendental
Functions,''\/} Vol.~II, McGraw-Hill, New York, 1953.

\item {{\bf [6]}} A. Erd\'elyi et al., {\it ``Tables of Integral
Transforms,''\/} Vol.~II, McGraw-Hill, New York, 1954.

\item {{\bf [7]}} J. Garc\'\i a Cuerva and J. L. Rubio de Francia,
{\it ``Weighted Norm Inequalities and Related Topics,''\/} North-Holland,
Amsterdam, 1985.

\item {{\bf [8]}} J. J. Guadalupe, M. P\'erez, F. J. Ruiz and J. L. Varona,
{\it Mean and weak convergence of Fourier-Bessel series}, J. Math. Anal. Appl.
(to appear).

\item {{\bf [9]}} J. J. Guadalupe, M. P\'erez and J. L. Varona, {\it Weak
behaviour of Fourier-Jacobi series}, J. Approx. Theory {\bf 61} (1990),
222--238.

\item {{\bf [10]}} J. J. Guadalupe, M. P\'erez and J. L. Varona, {\it Mean and
weak convergence of some orthogonal Fourier expansions by using $A_p$ theory},
in Orthogonal Polynomials and Their Applications (Proc. Int. Congr.,
Laredo/Spain 1987, Ed. J. Vinuesa), 161--169, Lect. Notes Pure Appl. Math., 117,
Dekker, New York, 1989.

\item {{\bf [11]}} R. Hunt, B. Muckenhoupt and R. Wheeden, {\it Weighted norm
inequalities for the conjugate function and the Hilbert transform},
Trans. Amer. Math. Soc. {\bf 176} (1973), 227--251.

\item {{\bf [12]}} J. Newman and W. Rudin, {\it Mean convergence of orthogonal
series}, Proc. Amer. Math. Soc. {\bf 3} (1952), 219--222.

\item {{\bf [13]}} H. Pollard, {\it The mean convergence of orthogonal
series I}, Trans. Amer. Math. Soc. {\bf 62} (1947), 387--403.

\item {{\bf [14]}} E. M. Stein and G. Weiss, {\it ``Introduction to Fourier
Analysis on Euclidean Spaces,''\/} Princeton Univ. Press, Princeton, NJ, 1975.

\item {{\bf [15]}} G. N. Watson, {\it ``A Treatise on the Theory of Bessel
Functions,''\/} Cambridge Univ. Press, 1966.

\item {{\bf [16]}} J. E. Wilkins, Jr., {\it Neumann series of Bessel functions},
Trans. Amer. Math. Soc. {\bf 64} (1948), 359--385.

\item {{\bf [17]}} J. E. Wilkins, Jr., {\it Neumann series of Bessel functions.
II}, Trans. Amer. Math. Soc. {\bf 69} (1950), 55--65.

\bigskip


% ------------------------------------------------


\noindent{\sl Juan L. Varona}\par
\noindent{\sl Departamento de Matem\'atica Aplicada}\par
\noindent{\sl Colegio Universitario de La Rioja (Universidad de Zaragoza)}\par
\noindent{\sl Obispo Bustamante, 3}\par
\noindent{\sl 26001 Logro\~no, SPAIN}




\bye
